\chapter{Revy}

\TKETs Revy er en stolt og lang tradition. Den første revy blev opført i maj 1958, og det har været en fast tradition med en årlig revy siden julerevyen 1962. I det sidste lange stykke tid har der været to årlige revyer: en julerevy med cancan-dans fra mandebest, og en maj-revy med en formtale skrevet af ceremonimesteren.

\section*{Roller}

For at lave en revy skal en masse mennesker lægge et masse arbejde. Her bringer vi lidt om de forskellige roller der indgår i revyen.

\subsection{Destruktøren}
I fordums var det \CERM{s} opgave at planlægge og afvikle revyerne, men siden 2014 har det været Destruktørens opgave at styre revyen med hård hånd. Destruktøren bliver, modsat andre ansvarsposter i foreningen faktisk \emph{valgt} af revyens medlemmer. Populært siges det at ''Destruktøren har Aben", hvorfor der også følger en lille abe-figur med rollen.

\subsection{Revyister}
Det er revyisternes ansvar at skrive samt fremføre revyen. I nyere tid involverer det et skrive/fremlægningsmøde hver uge, samt selvfølgelig øveweekender op til forestillingerne.

\subsection{Bandet\texttrademark}
\TKETs fantastiske huskortester Bandet\texttrademark \;spiller godt. Derudover holder de to årlige koncerter hhv. en jule- og maj koncert, hvor de reklamerer med, at der vil være sketches i pauserne mellem sangene. Revyen er lykkelig og beæret over at disse koncerter sammenfalder med revyen hvert år.

\subsection{Rekvitøser}
Til en revy er der brug for en forfærdelig masse rekvisitter. Revyen har en herskare af dygtige rekvitøser der kan kræerere alt fra en replika af klokken (2023 Maj) til en menneske-størrelse weber-grill kostume (2024 Jul). At dette kan lade sig gøre, er absurd impoenerende. Om muligt endnu mere imponerende er, at stort set alt sammen kan laves a pap og valgplakater\footnote{Næsten. Engang imellem er en rekvitøs ferm med en symaskine og kan lave eksempelvis en majs-hat (2021 Maj)}! Det er også rekvitøserne der hver revy maler det flotte bagtæppe.

Siden julerevyen 2024 har \CERM afviklet en kvit-auktion og solgt de sejeste kvitter fra revyen. Profitten kan derved gå til nye og sejere kvitter, og gæsterne kan få et minde med hjem. Dette har vist sig at være en meget indbringe forretning -- senest 3297 kroner til majrevyen 2026.

\subsection{Teknik}
Revyen har både lys og lyd. Det er teknik der står for at lyse på revyisterne, afspille de mange lydeffekter, samt sørge for at man kan høre både revyister og band i hele lokalet. Derudover filmer de også en revydeo, som revyholdet bruger til at evaluere processen.

\subsection{Tæppefører}
Det er svært at se revyen hvis tæppet er trukket for hele tiden. Heldigvis har revyen en tæppefører der kan stå for at trække det fra og lukke det på alle de skøre tidspunkter og i alle de skøre tempi revyisterne kan finde på.


\section*{Eksempler på indhold}
Som en smagsprøve på hvad der kan være med i en revy bringer vi her to fisk og en formtale. Vi ville gerne også have bragt et eksempel på cancan, men det var ikke nemt på skrift.

% Pepsi Twist

\begin{Sketch}{Pepsi Twist}
	
	\begin{Regi}
		\textbf{P1} og \textbf{P2} står ved cafébordet på forscenen. \textbf{T} kommer ind.
	\end{Regi}
	
	\begin{Replik}[T]
		Kan jeg tage imod bestillinger for drikkevarer?
	\end{Replik}
	
	\begin{Replik}[P1]
		Ja. Jeg kunne godt tænke mig en dansvand.
	\end{Replik}
	
	\begin{Replik}[P2]
		Og jeg vil gerne have ne Cola.
	\end{Replik}
	
	\begin{Replik}[T]
		Så gerne. Jeg er straks tilbage.
	\end{Replik}
	
	\begin{Regi}
		\textbf{T} går ud. \textbf{P1} og \textbf{P2} kan evt. snakke videre som om de var midt i en samtale. Det må gerne være obskurt. Eksempler kunnne være:
	\end{Regi}
	
	\begin{Replik}[P1]
		Så\dots Har du været ude og rejse for nyligt?
	\end{Replik}
	
	\begin{Replik}[P2]
		Jeg drikker ikke. Jeg er mormor.
	\end{Replik}
	
	\begin{Regi}
		Eller
	\end{Regi}
	
	\begin{Replik}[P1]
		Så\dots Har du set den nye Jumanji?
	\end{Replik}
	
	\begin{Replik}[P2]
		Jeg tror ikke på Jack Black.
	\end{Replik}
	
	\begin{Regi}
		\textbf{T} kommer ind igen efter kort tid med flaskerne.
	\end{Regi}
	
	\begin{Replik}[T]
		Så var der en danskvand til dig. Og en Cola til dig. Vi har ikke Coca Cola, så jeg gik ud fra at Pepsi var okay. (\lyd{Dramatisk musik})
	\end{Replik}
	
	\begin{Replik}[S]
		\TKET præsenterer: Pepsi Twist!
	\end{Replik}
	
\end{Sketch}

\begin{Footer}
	Bjarke Bilenberg Demant
\end{Footer}

% Platter Jokes

\begin{Sketch}{Platter Jokes}

\begin{Regi}
	K og Å kommer ind på scenen i festligt humør med en dyb tallerken hver. De kigger på hinanden og klinker deres tallerkener sammen som en skål.
\end{Regi}

\begin{Replik}[S]
	\TKET præsenterer: ... Ej, det er skidt.
\end{Replik}

\begin{Replik}[K]
	Jo, kom nu \regi{Henvendt til S}
\end{Replik}


\begin{Replik}[S]
	Okay, okay. \TKET præsenterer: Skål \regi{Noget uengageret}
\end{Replik}
% Sådan skriver du en replik - her inkl en regibemærkning 
\begin{Replik}[K]
	Hvad gik det ud på? Publikum var heller ikke så vilde med den, som jeg troede
\end{Replik}


\begin{Replik}[Å]
	Jeg ved ikke helt hvad der skete. Tror du, det er fordi, de ikke forstår den?
\end{Replik}

\begin{Replik}[K]
	\regi{Henvendt til publikum} Altså ser I, det er fordi, at vi har to dybe tallerkener med os her
\end{Replik}

\begin{Regi}
	K og Å gestikulerer med deres tallerkener. 
\end{Regi}

\begin{Replik}[Å]
	Ja, og altså, her i Danmark gør vi tit det, at vi ligesom slår vores glas sammen og siger skål.
\end{Replik}

\begin{Regi}
	K og Å klinker deres tallerkener igen og siger skål.
\end{Regi}

\begin{Replik}[Å]
	Og et andet ord for dybe tallerkener er skål, så det er jo faktisk det samme ord, men med to betydninger.
\end{Replik}

\begin{Replik}[K]
	Det er et slags ordspil, I ved
\end{Replik}

\begin{Replik}[Å]
	Ja, eller et spil på ord, om man vil \regi{Å blinker overdrevet}
\end{Replik}

\begin{Replik}[K]
	Faktisk, hvis vi går tilbage til ordets oprindelse, kommer det fra det oldnordiske skól, som betyder hovedskalle. Det er en reference til, at den kræver lidt tankekraft at forstå den her joke.
\end{Replik}
\begin{Replik}[Å]
	Den er faktisk sjovere, end man lige skulle tro.
\end{Replik}

\begin{Regi}
	K og Å kigger meget entusiastisk ud på publikum, men ser mere og mere forvirret ud og kigger til sidst på hinanden.
\end{Regi}

\begin{Replik}[K]
	Hva, den er egentligt ikke så god, er den det?
\end{Replik}

\begin{Replik}[Å]
	Nej, det er den egentligt ikke. Kvaliteten er virkeligt dalet i revyen i år, synes jeg. Men det kan jo heller ikke blive ved med at være godt.
\end{Replik}

\begin{Replik}[K]
	Nej, du har ret. Det er jo ikke hver gang, at man kan opfinde den dybe tallerken.
\end{Replik}

\begin{Regi}
	K og Å laver overdrevet bevægelse ud mod publikum á la Looney Tunes og Bandet$^{\textbf{TM}}$ spiller Ba-Dum-Tss
\end{Regi}

\begin{Regi}
	K og Å går af scenen.
\end{Regi}

\end{Sketch}

\begin{Footer}
Tobias Lindholdt Johannesen (\TKprefix{2025}\CERM)
\end{Footer}

\subsection{Formtale, Generalprøven Maj 2024}

\noindent Kære medlemmer og gæster, velkommen til midten af forestillingen. Og velkommen til mig, \TKETS formand.

\noindent I løbet af alle mine mange knap et år som FORM, er jeg blevet fortalt, at nogle personer har svært ved at forstå mig. Dette har jeg dog svært ved at forestille mig. Hvad er svært at forstå ved sætningen ``på bordet eller på tavlen?'' Eller sætningen ``jeg har ret.'' Eller sætningen ``\(\neg(\) jeg ville ikke hade, hvis der ikke gjalt, at jeg ønskede, at der ikke var en dag og/eller en nat, i løbet af hvilken jeg tænkte på andet end, at NF ikke er klog \()\).''

\noindent Men for at hjælpe dem, som ikke kan forstå mig, har jeg lært et universelt sprog: Nemlig latin! Fremadrettet vil alle torsdagsmøder foregå på latin. For at give en forsmag, får I her et eksempel:

[\textbf{REGI: synges til melodien ``I en kælder sort som kul.''}]\\
\noindent Lorem ipsum dolor sit amet, consectetur adipiscing elit, sed do eiusmod tempor incididunt ut labore et dolore magna aliqua. Cantabo nunc ex corde meo de \TKET. de osculo somnio NF; osculum tam pulchrum quam SEKR. Spero PR osculum pingit. Et volo pro KASS solvere pro cornibus pro \CERM et \VC. Ego ipsa non volo INKA verberare.

\noindent Det føltes godt endeligt at blive forstået af alle. Hvis I har brug for endnu et eksempel, så spørg mig til festen; så fremviser jeg en valgfri sketch fra revyen på latin, hvor jeg spiller alle roller. Desuden mener jeg, at \TKET vil bestå.

\begin{Footer}
	Zacharias Niclasen (\TKprefix{2024}\CERM)
\end{Footer}